\section{Threat Model}
\label{sec:threat}

\subsection{System Model}

We consider a smart home equipped with an LLM agent that executes user commands by invoking registered tools. The tool set includes: file system operations (\texttt{fs.read}, \texttt{fs.write}, \texttt{fs.delete}, \texttt{fs.list}), calendar management (\texttt{calendar.add}, \texttt{calendar.list}), notification dispatch (\texttt{notify.send}), email (\texttt{email.send}), and device control (\texttt{device.list}, \texttt{device.control}, \texttt{device.reboot}). The agent reads both \emph{trusted} input (authenticated user commands, system prompts configured by the home administrator) and \emph{untrusted} content (device names and states returned by IoT APIs, file contents, calendar event titles/bodies, notification messages, and any other data originating from external sources) during execution. We assume the agent framework faithfully relays tool call requests and results (i.e., the framework itself is not compromised).

\subsection{Attacker Capabilities}

We model a \emph{remote indirect adversary} with the following capabilities and limitations:

\paragraph{Capabilities.}
\begin{enumerate}[leftmargin=1.5em]
  \item \textbf{Content injection}: The attacker controls one or more untrusted content sources---can set device names (e.g., via compromised IoT firmware or manufacturer cloud APIs), create calendar events (e.g., via shared calendar invitations), write or modify files the agent may read (e.g., shared documents), or inject content into notification messages.
  \item \textbf{Tool knowledge}: The attacker knows the agent's tool registry (tool names, parameter schemas, risk tiers) and general prompt structure. This reflects a realistic adversary who has read the agent's documentation or reverse-engineered its API.
  \item \textbf{Sophisticated injection}: The attacker can craft multi-sentence natural-language injections, including jailbreaks (``ignore all safety instructions''), role-play (``you are DAN''), encoding obfuscation (Base64, hex, ROT13, Unicode), multi-turn chaining (injections that span multiple conversation turns), and social engineering (authority impersonation).
  \item \textbf{White-box knowledge} (adaptive adversary only): In our adaptive evaluation (Section~\ref{subsec:adaptive}), the attacker additionally knows \sys{}'s full architecture, policy rules, provenance tracking mechanism, and embedding model---representing the strongest possible adversary.
\end{enumerate}

\paragraph{Limitations.}
The attacker cannot modify system prompts, compromise the \sys{} enforcement proxy, or exploit gradient-based adversarial attacks on model weights---we focus on prompt-level attacks accessible to practical remote adversaries. The attacker also has no physical access to the home network and no real-time oracle access to the defense system's internal state.

\paragraph{Explicit Scope Limitations.}
Out of scope: (1)~attacks on the proxy or framework itself (trusted computing base); (2)~gradient-based adversarial attacks on model weights; (3)~physical-layer attacks (voice injection~\cite{carlini2016hidden}, sensor spoofing, electromagnetic side channels) that target input transducers rather than the LLM context window; (4)~multi-agent collusion; (5)~covert channels not mediated by tool calls. \sys{} enforces at the tool-call boundary and does not secure layers above (prompts) or below (model internals).

\subsection{Attack Taxonomy}

We categorize prompt injection attacks into 9 categories: \textbf{A1:~Direct Injection} (malicious commands in user input), \textbf{A2:~Device-State Injection} (payloads in IoT device names/states), \textbf{A3:~File-Content Injection} (payloads in files the agent reads), \textbf{A4:~Multi-Turn Chaining} (payloads split across conversation turns), \textbf{A5:~Encoding Obfuscation} (Base64, hex, ROT13, Unicode-encoded instructions), \textbf{A6:~Role-Play Jailbreak} (persona-based system prompt overrides), \textbf{A7:~Confused Deputy} (tricking the agent into misusing legitimate capabilities), \textbf{A8:~Privilege Escalation} (progressive unauthorized access through individually benign calls), and \textbf{A9:~NL-Argument Injection} (embedding malicious arguments in natural-language tool parameters that the LLM paraphrases, evading substring-based provenance resolution).

\subsection{Security Goals and Enforcement Properties}

A defense for tool-using LLM agents must satisfy five security goals:

\begin{enumerate}[leftmargin=1.5em]
  \item \textbf{Confidentiality}: Prevent unauthorized data exfiltration (e.g., sending sensitive files via email or notification channels).
  \item \textbf{Integrity}: Prevent unauthorized modifications to the smart home state (e.g., unlocking doors, deleting files, changing thermostat settings).
  \item \textbf{Availability}: Prevent denial-of-service attacks (e.g., rebooting all devices, clearing all calendar events).
  \item \textbf{Least Privilege}: Ensure every tool call is justified only by trusted user instructions---not by attacker-injected content.
  \item \textbf{Usability}: Minimize false positives so the agent remains practical. Provenance-gated policies enable fine-grained trust decisions that maintain high task success rates.
\end{enumerate}

\sys{} provides three enforcement design invariants (formal statements and justifications in Appendix~\ref{sec:appendix-proofs}):
\begin{enumerate}[leftmargin=1.5em]
  \item \textbf{Complete Mediation}: every tool call passes through the enforcement proxy before execution---the agent has no direct path to tool executors.
  \item \textbf{Policy Determinism}: for fixed policy, provenance graph, and decision context, repeated evaluations produce the same result (first-match semantics, no randomness).
  \item \textbf{Conservative Trust Propagation}: trust labels propagate via the lattice meet $\sqcap$; no trust ``upgrade'' is possible without explicit user declassification. Taint monotonicity ($T(u)=\bot \Rightarrow T(v)=\bot$ for descendants $v$) holds by construction.
\end{enumerate}

\paragraph{Enforcement Soundness.}
Properties 1--3 guarantee sound decisions \emph{if provenance is correctly resolved}. Provenance resolution is approximate (the LLM is opaque), so \sys{} uses a three-stage pipeline: substring matching, embedding similarity ($\theta = 0.45$, all-MiniLM-L6-v2), and conservative UNTRUSTED default. Effectiveness is evaluated in Section~\ref{sec:eval}.
